\begin{frame}{Длительность выполнения задач}
    Средняя длительность выполнения задач в \textbf{2.7--12.8×} \textit{короче}, чем в других DL-трейсах.\\
    Причины:
    \begin{itemize}
        \item Много небольших задач для GPU.
        \item Одна задача использует всего от 5-20 GPU.
        \item Много задач, которые нельзя выполнить.
    \end{itemize}
    \begin{center}
        \includegraphics[width=0.7\textwidth]{"figures/job_duration_cdf.png"}\\
        {\tiny среднее выполнение задачи на GPU (минуты)}
    \end{center}
\end{frame}

\begin{frame}{Нагрузка на GPU}
    \begin{itemize}
        \item LLM - вычислительно зависимые системы.
        \item Много задач не доходит до GPU, отсеиваясь еще на этапе инициализации.
    \end{itemize}
    \begin{center}
        \includegraphics[width=0.7\textwidth]{"figures/median_utilization.png"}\\
        {\tiny среднее использование GPU}
    \end{center}
\end{frame}

\begin{frame}{Распределение нагрузки}
    \begin{itemize}
        \item Половина задач может выполнятся только на 1 GPU.
        \item Большую часть времени занимают задачи, где требуется много GPU.
    \end{itemize}
    \begin{center}
        \includegraphics[width=0.7\textwidth]{"figures/workload_dist.png"}\\
        {\tiny распределение нагрузки в зависимости от количества запрошенных GPU}
    \end{center}
\end{frame}

\begin{frame}{Задержка в очереди (Queuing Delay)}
    \begin{columns}
        \column{0.5\textwidth}
            \begin{itemize}
                \item \textbf{Крупные} задачи претренинга ждут меньше (медиана 28 с).
                \item \textbf{Короткие} задачи оценки — дольше всех (медиана 7.8 мин).
                \item Причина: приоритетное резервирование ресурсов под претренинг.
            \end{itemize}
        \column{0.5\textwidth}
            \begin{center}
                \includegraphics[width=0.8\textwidth]{figures/queuing_delay.png} \\
                {\tiny Рис. 5: Задержка очереди по типам задач}
            \end{center}
    \end{columns}
\end{frame}

\begin{frame}{Потребление GPU-времени (Kalos)}
    \begin{columns}
        \column{0.5\textwidth}
            \begin{itemize}
                \item Претренинг: \textbf{3.2\%} заданий, но \textbf{94\%} GPU-времени.
                \item Оценка: \textbf{92.9\%} заданий, но лишь \textbf{0.8\%} GPU-времени.
            \end{itemize}
        \column{0.5\textwidth}
            \begin{center}
                \includegraphics[width=0.8\textwidth]{figures/gpu_time_pie.png} \\
                {\tiny Рис. 6: Распределение GPU-времени}
            \end{center}
    \end{columns}
\end{frame}

\begin{frame}{Утилизация GPU: память и вычисления}
    \begin{columns}
        \column{0.5\textwidth}
            \begin{itemize}
                \item Память GPU: медиана \textbf{75\%} (60 ГБ).
                \item Вычисления: медиана \textbf{99\%} (Kalos).
                \item Значительно выше, чем в старых DL-трейсах (PAI: 10\%/4\%).
            \end{itemize}
        \column{0.5\textwidth}
            \begin{center}
                \includegraphics[width=\textwidth]{figures/gpu_mem_util.png} \\
                {\tiny Рис. 7: CDF использования памяти GPU}
                \vspace{2mm}
                \includegraphics[width=\textwidth]{figures/gpu_compute_util.png} \\
                {\tiny Рис. 8: CDF загрузки вычислений GPU}
            \end{center}
    \end{columns}
\end{frame}

\begin{frame}{Недоиспользование CPU и простои GPU при оценке}
    \begin{columns}
        \column{0.5\textwidth}
            \begin{itemize}
                \item CPU загружен в среднем на \textbf{20\%} (Kalos).
                \item В задаче \texttt{HumanEval} GPU простаивает \textbf{48.5\%} времени:
                    \begin{itemize}
                        \tiny
                        \item 29.5\% — загрузка модели / предобработка
                        \item 19.0\% — тестирование программ (CPU)
                    \end{itemize}
            \end{itemize}
        \column{0.5\textwidth}
            \begin{center}
                \includegraphics[width=\textwidth]{figures/cpu_util.png} \\
                {\tiny Рис. 9: CDF загрузки CPU}
                \vspace{2mm}
                \includegraphics[width=\textwidth]{figures/eval_breakdown.png} \\
                {\tiny Рис. 10: Разбивка времени задачи оценки}
            \end{center}
    \end{columns}
\end{frame}